%---------------------------------------------------------------------------%
%->> Frontmatter
%---------------------------------------------------------------------------%
%-
%-> 生成封面
%-
\maketitle% 生成中文封面
\MAKETITLE% 生成英文封面
%-
%-> 作者声明
%-
%-
%-> 中文摘要
%-
\let\cleardoublepage\clearpage  % 关键：取消章节强制从奇数页开始
\intobmk\chapter*{摘\quad 要}% 显示在书签但不显示在目录
\setcounter{page}{1}% 开始页码
\pagenumbering{Roman}% 页码符号

本项目围绕 PostgreSQL master 分支（2022–2025）的提交历史，构建“提交行为统计—Bug 修复识别—静态结构演化—动态执行追踪—形式化验证”的多层分析链路，系统刻画 Bug 修复行为规律与代码演化特征，并对关键修复进行可验证性证明。在提交层面，基于规则化标注与时间/作者/模块统计，发现修复节奏与发布周期相关，高风险模块（storage/replication/WAL 等）长期高占比，且核心开发者承担关键模块的修复工作。在静态结构层面，Bug 修复呈现出一致的结构性演化模式：前置条件显式化、控制流分裂、错误路径显式化与条件语义增强等，体现防御式编程的系统性强化。在动态分析层面，通过 PySnooper 追踪 PostgreSQL Python 工具脚本，定位出配置同步脚本的“静默逻辑失效”，并通过回归测试验证修复有效性，建立了从异常发现到验证的闭环流程。在形式化验证层面，使用 Z3 对该修复进行约束建模，证明原逻辑存在反例（SAT），修复后逻辑在规约域内无反例（UNSAT），从而实现“证明修复有效”的强结论。整体而言，本研究以统计分析、结构分析、动态追踪与形式化建模互证，为理解大型系统的 Bug 修复规律与提升修复可信度提供了可复用的方法框架。

\keywords{PostgreSQL，提交历史分析，Bug 修复识别，代码结构演化，动态追踪，形式化验证}% 中文关键词
%-
%-> 英文摘要
%-
\intobmk\chapter*{Abstract}% 显示在书签但不显示在目录

This thesis builds a multi-layer analysis pipeline over the PostgreSQL master branch commit history (2022--2025), integrating commit-level statistics, bug-fix identification, static structural evolution analysis, dynamic execution tracing, and formal verification. At the commit level, using rule-based labeling and statistics by time, author, and subsystem, we observe that fix activity correlates with release cycles; high-risk subsystems (e.g., storage, replication, and WAL) consistently account for a large share of fixes, and core developers contribute disproportionately to critical-module maintenance. At the static level, bug fixes exhibit recurring structural evolution patterns, including explicit preconditions, control-flow splitting, explicit error paths, and enriched conditional semantics, reflecting systematic strengthening of defensive programming. At the dynamic level, we trace PostgreSQL Python utility scripts with PySnooper to pinpoint a silent logic failure in a configuration synchronization script, validate the fix via regression tests, and close the loop from anomaly discovery to verification. Finally, we model the fix in Z3: the original logic admits counterexamples (SAT), whereas the repaired logic has no counterexample within the specification domain (UNSAT), yielding a strong, proof-backed validity claim. Overall, the proposed framework triangulates statistical analysis, structural analysis, dynamic tracing, and formal modeling to provide reusable methods for understanding bug-fix behaviors in large systems and improving the trustworthiness of fixes.


\KEYWORDS{PostgreSQL, commit history mining, bug-fix identification, structural evolution, dynamic tracing, formal verification}% 英文关键词
%---------------------------------------------------------------------------%
